\section{The family of order \texorpdfstring{$6k+6$}{6k+6}}\label{sec:family6}
The same layer construction also applies to a twelve-vertex seed at strict
level one. This example uses five auxiliary types, and its positive
certificates can be written out entirely as folds inside successive
neighbourhoods. It illustrates how the three finite hypotheses are checked
when no recursive certificate file is needed.
We first state the finite assertions and their consequences; the reductions
and clique tables used in the proof are displayed in
Section~\ref{6k6:app:certificates}.

\subsection{The seed and the resulting family}\label{6k6:sec:seed}
Here five auxiliary types give six new vertices per layer. Let
\(V(A)=\{a_0,\ldots,a_{11}\}\), with the adjacency list in
Table~\ref{6k6:tab:Aadj}. The list is symmetric, has no loops, and gives
\(\card{E(A)}=47\). The profiles are given in Table~\ref{6k6:tab:profiles};
their cardinalities sum to \(48\). The type graph is
\begin{equation}\label{6k6:eq:R}
 E(R)=\{02,03,04,12,14,23,34\},
\end{equation}
where \(jt\) abbreviates the unordered pair \(\{j,t\}\).
There are no other edges of \(R\).
\begin{table}[H]
\centering\small\setlength{\tabcolsep}{4pt}
\caption{The complete adjacency list of \(A\). An integer \(i\) denotes \(a_i\).}\label{6k6:tab:Aadj}
\begin{tabular}{c l}
\toprule
$v$ & $N_A(v)$\\
\midrule

0 & \(\{1,3,4,5,6,7,9,10\}\)\\
1 & \(\{0,2,3,4,5,6,7,8,10\}\)\\
2 & \(\{1,4,5,6,8,9,11\}\)\\
3 & \(\{0,1,4,6,7,8,9,11\}\)\\
4 & \(\{0,1,2,3,5,8,10,11\}\)\\
5 & \(\{0,1,2,4,6,9,10,11\}\)\\
6 & \(\{0,1,2,3,5,8,9\}\)\\
7 & \(\{0,1,3,9,10,11\}\)\\
8 & \(\{1,2,3,4,6,9,10,11\}\)\\
9 & \(\{0,2,3,5,6,7,8,10,11\}\)\\
10 & \(\{0,1,4,5,7,8,9,11\}\)\\
11 & \(\{2,3,4,5,7,8,9,10\}\)\\
\bottomrule
\end{tabular}
\end{table}

\begin{table}[H]
\centering\small\setlength{\tabcolsep}{4pt}
\caption{The five profiles. Each displayed integer \(i\) denotes \(a_i\).}\label{6k6:tab:profiles}
\begin{tabular}{c l}\toprule
$j$ & Indices of the vertices in $P_j$\\\midrule
0 & \(1,2,3,4,5,6,7,8,9,11\)\\
1 & \(1,3,4,6,7,8,10,11\)\\
2 & \(0,1,3,4,5,6,7,9,10,11\)\\
3 & \(1,2,3,5,6,7,8,9,10,11\)\\
4 & \(0,1,2,3,5,7,8,9,10,11\)\\
\bottomrule\end{tabular}
\end{table}

The graph \(B\) consists of \(A\), the centre \(p\), and the five
auxiliaries \(x_0,\ldots,x_4\). Its centre is adjacent to every vertex
of \(A\) and to no auxiliary; the neighbours of \(x_j\) in \(A\)
are precisely \(P_j\), and its auxiliary adjacencies are those of \(R\).
Thus
\[
 \card{V(B)}=18,\qquad
 \card{E(B)}=47+12+48+7=114.
\]
In the finite certificates, \(a_i\) is abbreviated to \(i\), and
\(p=12\), \(x_j=13+j\). In particular, vertex \(15\) denotes \(x_2\).
\begin{proposition}\label{6k6:prop:seed}
These seed data satisfy
\[
 A\in\D_1\setminus\D_0,\qquad B-p\in\D_1,\qquad
 q(B[N_B(v)])=
 \begin{cases}
 2,&v=a_7\text{ or }v=x_2,\\
 1,&v\ne p,a_7,x_2
 \end{cases}
 \quad(v\ne p).
\]
\end{proposition}
\begin{proof}
Table~\ref{6k6:tab:Ano0} gives, for every \(v\in V(A)\), the degree
\(d_v=\card{N_A(v)}\) and the maximum degree in \(A[N_A(v)]\).
In each row that maximum is strictly smaller than \(d_v-1\).
No open neighbourhood graph is a cone; Lemma~\ref{lem:first}
therefore gives \(A\notin\D_0\).

Tables~\ref{6k6:tab:Apos} and~\ref{6k6:tab:Bpos} specify successive outer
deletions in \(A\) and \(B-p\), respectively. Each row lists the
exact neighbourhood set in the current outer graph, the folds
within that neighbourhood in the stated order, and an apex of
the resulting neighbourhood. All fourteen required closed
neighbourhood inclusions are displayed in Table~\ref{6k6:tab:inclusiones}.
They are obtained by restricting the adjacency data of
Tables~\ref{6k6:tab:Aadj} and~\ref{6k6:tab:profiles} and~\eqref{6k6:eq:R} to the
current set. After the outer deletions, the remaining vertex sets
are cones with the apices specified in Section~\ref{6k6:app:certificates}.
Each inner reduction is consequently a level-\(0\) reduction, and
each outer deletion is permitted at level \(1\).
The final cones also lie in \(\D_1\), by Lemma~\ref{lem:inclusions}.
This proves \(A\in\D_1\) and \(B-p\in\D_1\).

Table~\ref{6k6:tab:Bq} gives all clique coefficients of the seventeen
non-central neighbourhood graphs. These coefficients follow by
the coefficient recurrence in Lemma~\ref{lem:qrec}, with
\(c_0(H)=1\) and \(c_t(\varnothing)=0\) for \(t\geq1\).
Their alternating sums are \(2\) in rows \(7\) and \(15\), and
\(1\) in every other row. This proves the last assertion.
\end{proof}

For this seed, the graphs of Section~\ref{sec:construction} satisfy
\begin{equation}\label{6k6:eq:concrete-order}
 \card{V(G_m)}=12+6m,\qquad
 \card{E(G_m)}=47+52m+15m^2 .
\end{equation}
Indeed, the seven disjoint classes of edges give
\[
\begin{aligned}
 \card{E(G_m)}
 &=47+12m+48m+\binom m2+5\binom m2+7m^2+5m(m-1)\\
 &=47+52m+15m^2.
\end{aligned}
\]
The terms count, in order, edges within \(A\), between \(A\) and
centres, between \(A\) and auxiliaries, between centres, within
each auxiliary type, between distinct types, and between centres
and auxiliaries.

\begin{theorem}\label{6k6:thm:family}
For every $k\geq1$, put $J_k=G_{k-1}$. Then $J_k\in\D_k\setminus\D_{k-1}$ and $|V(J_k)|=6k+6$.
\end{theorem}
\begin{proof}
Proposition~\ref{6k6:prop:seed} supplies the three hypotheses of
Theorem~\ref{thm:template}, with \(s=1\) and \(d=5\). Consequently
\[
 G_m\in\D_{m+1}\setminus\D_m\qquad(m\geq0).
\]
Taking \(m=k-1\) gives \(J_k\in\D_k\setminus\D_{k-1}\) and
\(\card{V(J_k)}=12+6(k-1)=6k+6\).
\end{proof}

\subsection{The finite reductions and obstructions}\label{6k6:app:certificates}
The certificates below use only the adjacency list, the profiles,
and the type edges specified in Section~\ref{6k6:sec:seed}.

\subsubsection{Excluding a first fold in \texorpdfstring{\(A\)}{A}}
For \(v\in V(A)\), put \(d_v=\card{N_A(v)}\), and write
\(\Delta(Q)\) for the maximum degree in a non-empty graph \(Q\).
The entries of Table~\ref{6k6:tab:Ano0} satisfy
\(\Delta(A[N_A(v)])<d_v-1\), so none of these neighbourhoods
has an apex.
\begin{table}[H]
\centering\small\setlength{\tabcolsep}{4pt}
\caption{Every open neighbourhood graph of \(A\) has no universal vertex.}\label{6k6:tab:Ano0}
\begin{tabular}{c r r r}
\toprule
$v$ & $d_v=\card{N_A(v)}$ & $d_v-1$ & $\Delta(A[N_A(v)])$\\
\midrule

0 & 8 & 7 & 6\\
1 & 9 & 8 & 6\\
2 & 7 & 6 & 5\\
3 & 8 & 7 & 5\\
4 & 8 & 7 & 6\\
5 & 8 & 7 & 5\\
6 & 7 & 6 & 5\\
7 & 6 & 5 & 4\\
8 & 8 & 7 & 5\\
9 & 9 & 8 & 6\\
10 & 8 & 7 & 5\\
11 & 8 & 7 & 6\\
\bottomrule
\end{tabular}
\end{table}


\subsubsection{The two level-1 reductions}
In each row of Tables~\ref{6k6:tab:Apos} and~\ref{6k6:tab:Bpos}, \(v\) is
deleted from the outer graph after all preceding outer deletions.
The set \(S\) is its neighbourhood set at that moment.
The arrows are folds inside the graph induced by \(S\).
After each arrow its source vertex is removed before the next
inclusion is evaluated. The listed apex is universal after
those inner folds. A row with no folds already has a conical
neighbourhood graph.

For Table~\ref{6k6:tab:Apos}, the final outer vertex set is
\(\{5,6,8,9,10,11\}\), with apex \(9\).
For Table~\ref{6k6:tab:Bpos}, it is
\(\{3,7,8,9,10,11,13,14,15,16,17\}\), with apex \(11\).
The centre \(p\) has been removed before Table~\ref{6k6:tab:Bpos}
begins.
\begin{table}[H]
\centering\small\setlength{\tabcolsep}{4pt}
\caption{A reduction certificate for \(A\in\D_1\).}\label{6k6:tab:Apos}
\begin{tabular}{c p{0.45\textwidth} p{0.22\textwidth} c}
\toprule
$v$ & Current neighbourhood set $S$ & Folds in order & Apex\\
\midrule

7 & \(\{0,1,3,9,10,11\}\) & \(1\to 0\) & 9\\
0 & \(\{1,3,4,5,6,9,10\}\) & \(4\to 1\), \(3\to 6\) & 5\\
1 & \(\{2,3,4,5,6,8,10\}\) & \(3\to 8\), \(6\to 2\) & 4\\
2 & \(\{4,5,6,8,9,11\}\) & \(4\to 11\) & 9\\
3 & \(\{4,6,8,9,11\}\) & none & 8\\
4 & \(\{5,8,10,11\}\) & none & 10\\
\bottomrule
\end{tabular}
\end{table}

\Needspace{16\baselineskip}
\begin{table}[H]
\centering\small\setlength{\tabcolsep}{4pt}
\caption{A reduction certificate for \(B-p\in\D_1\).}\label{6k6:tab:Bpos}
\begin{tabular}{c p{0.45\textwidth} p{0.22\textwidth} c}
\toprule
$v$ & Current neighbourhood set $S$ & Folds in order & Apex\\
\midrule

0 & \(\{1,3,4,5,6,7,9,10,15,17\}\) & \(4\to 1\), \(6\to 15\), \(3\to 7\) & 10\\
2 & \(\{1,4,5,6,8,9,11,13,16,17\}\) & none & 13\\
4 & \(\{1,3,5,8,10,11,13,14,15\}\) & \(5\to 15\), \(10\to 14\) & 3\\
1 & \(\{3,5,6,7,8,10,13,14,15,16,17\}\) & \(5\to 16\), \(6\to 3\), \(8\to 17\) & 7\\
5 & \(\{6,9,10,11,13,15,16,17\}\) & none & 9\\
6 & \(\{3,8,9,13,14,15,16\}\) & none & 3\\
\bottomrule
\end{tabular}
\end{table}

For example, the first row of Table~\ref{6k6:tab:Apos} starts with
\(S=\{0,1,3,9,10,11\}\). Within this neighbourhood,
\[
 N_{A[S]}[1]=\{0,1,3,10\}\subseteq
 N_{A[S]}[0]=\{0,1,3,9,10\}.
\]
The fold \(1\to0\) leaves a cone with apex \(9\), so the outer
vertex \(7\) can be deleted at level \(1\).
The next row is evaluated in \(A-7\).
Table~\ref{6k6:tab:inclusiones} provides the corresponding sets for
every fold in both certificates. Its third column is contained
in its fourth in each row; both columns are the closed
neighbourhoods in the current inner graph.
\begin{table}[H]
\centering\small\setlength{\tabcolsep}{4pt}
\caption{Closed neighbourhood inclusions for all fourteen folds. Here \(T\) is the neighbourhood graph after the preceding folds in the same row of the corresponding certificate.}\label{6k6:tab:inclusiones}
\begin{tabular}{c c p{0.31\textwidth} p{0.36\textwidth}}
\toprule
Graph; outer $v$ & Fold $u\to t$ & $N_T[u]$ & $N_T[t]$\\
\midrule

$A$; 7 & \(1\to 0\) & \(\{0,1,3,10\}\) & \(\{0,1,3,9,10\}\)\\
$A$; 0 & \(4\to 1\) & \(\{1,3,4,5,10\}\) & \(\{1,3,4,5,6,10\}\)\\
$A$; 0 & \(3\to 6\) & \(\{1,3,6,9\}\) & \(\{1,3,5,6,9\}\)\\
$A$; 1 & \(3\to 8\) & \(\{3,4,6,8\}\) & \(\{2,3,4,6,8,10\}\)\\
$A$; 1 & \(6\to 2\) & \(\{2,5,6,8\}\) & \(\{2,4,5,6,8\}\)\\
$A$; 2 & \(4\to 11\) & \(\{4,5,8,11\}\) & \(\{4,5,8,9,11\}\)\\
$B-p$; 0 & \(4\to 1\) & \(\{1,3,4,5,10,15\}\) & \(\{1,3,4,5,6,7,10,15,17\}\)\\
$B-p$; 0 & \(6\to 15\) & \(\{1,3,5,6,9,15\}\) & \(\{1,3,5,6,7,9,10,15\}\)\\
$B-p$; 0 & \(3\to 7\) & \(\{1,3,7,9,15,17\}\) & \(\{1,3,7,9,10,15,17\}\)\\
$B-p$; 4 & \(5\to 15\) & \(\{1,5,10,11,13,15\}\) & \(\{1,3,5,10,11,13,14,15\}\)\\
$B-p$; 4 & \(10\to 14\) & \(\{1,8,10,11,14,15\}\) & \(\{1,3,8,10,11,14,15\}\)\\
$B-p$; 1 & \(5\to 16\) & \(\{5,6,10,13,15,16,17\}\) & \(\{3,5,6,7,8,10,13,15,16,17\}\)\\
$B-p$; 1 & \(6\to 3\) & \(\{3,6,8,13,14,15,16\}\) & \(\{3,6,7,8,13,14,15,16,17\}\)\\
$B-p$; 1 & \(8\to 17\) & \(\{3,8,10,13,14,16,17\}\) & \(\{3,7,8,10,13,14,16,17\}\)\\
\bottomrule
\end{tabular}
\end{table}


\subsubsection{Clique coefficients}
Table~\ref{6k6:tab:Bq} supplies the coefficients \(c_0,\ldots,c_6\)
of every non-central neighbourhood graph in \(B\).
There are no cliques of greater cardinality in these graphs.
The coefficient recurrence~\eqref{eq:clique-coefficients}
computes every row from the adjacency data.
For instance, the alternating sums in rows \(0\) and \(7\) are
\[
 1-11+41-63+35-2=1,\qquad
 1-12+49-86+63-15+2=2.
\]
All seventeen sums are non-zero.
\begin{table}[H]
\centering\small\setlength{\tabcolsep}{4pt}
\caption{Clique counts for the seventeen non-central open neighbourhood graphs of \(B\). Coefficients \(\kappa_t\) with \(t>6\) are zero.}\label{6k6:tab:Bq}
\begin{tabular}{c r r r r r r r r}
\toprule
$v$ & $\kappa_0$ & $\kappa_1$ & $\kappa_2$ & $\kappa_3$ & $\kappa_4$ & $\kappa_5$ & $\kappa_6$ & $q$\\
\midrule

0 & 1 & 11 & 41 & 63 & 35 & 2 & 0 & 1\\
1 & 1 & 15 & 72 & 133 & 88 & 12 & 0 & 1\\
2 & 1 & 11 & 43 & 75 & 58 & 17 & 2 & 1\\
3 & 1 & 14 & 65 & 122 & 91 & 23 & 3 & 1\\
4 & 1 & 12 & 46 & 71 & 39 & 2 & 0 & 1\\
5 & 1 & 13 & 57 & 103 & 74 & 16 & 1 & 1\\
6 & 1 & 12 & 46 & 75 & 51 & 10 & 0 & 1\\
7 & 1 & 12 & 49 & 86 & 63 & 15 & 2 & 2\\
8 & 1 & 13 & 57 & 103 & 75 & 18 & 2 & 1\\
9 & 1 & 14 & 66 & 131 & 111 & 37 & 5 & 1\\
10 & 1 & 13 & 56 & 97 & 61 & 7 & 0 & 1\\
11 & 1 & 14 & 65 & 125 & 101 & 32 & 5 & 1\\
13 & 1 & 13 & 59 & 119 & 110 & 42 & 5 & 1\\
14 & 1 & 10 & 33 & 45 & 22 & 0 & 0 & 1\\
15 & 1 & 13 & 57 & 103 & 74 & 15 & 1 & 2\\
16 & 1 & 13 & 60 & 125 & 120 & 47 & 5 & 1\\
17 & 1 & 13 & 56 & 104 & 85 & 28 & 4 & 1\\
\bottomrule
\end{tabular}
\end{table}


\runcheck{6k6}
The output confirms twelve outer deletions, fourteen inner folds, twelve
rows excluding a first fold in the seed, and seventeen clique rows.
Each displayed neighbourhood and both sides of every fold inclusion are
recomputed from the adjacency lists. The calculation also verifies the
final cones in the two reductions.
